\thispagestyle{empty}
\begin{center}
\vspace*{26mm}
{\fontsize{26}{30}\selectfont\addfontfeatures{LetterSpace=10}SIDEREVS NVNCIVS}\\[4mm]
{\fontsize{16}{20}\selectfont\addfontfeatures{LetterSpace=8}THE STARRY MESSENGER}\\[10mm]
{\large Galilaeus Galilaeus · Galileo Galilei}\\[2mm]
{\small Venetiis, MDCX · Venice, 1610}\\[20mm]
{\color{edgray}\small Latin–English parallel edition}\\[6mm]
{\small edited by metachow}\\[1mm]
{\color{edgray}\footnotesize translated, proofread and typeset with Claude (Anthropic)}
\end{center}

\vfill
\noindent\begin{minipage}{\textwidth}
\small\color{edgray}\setlength{\parindent}{1.2em}\linespread{1.15}\selectfont
\textbf{Editorial note}\quad
The Latin follows A. Favaro's National Edition of Galileo's works, vol.\ III (1892), checked page by page against the first edition (Venice, 1610), in the Smithsonian Libraries copy digitised by the Internet Archive. The spelling is modernised: æ and œ are written out, long s is printed as s, u and v are distinguished, abbreviations are expanded, and classical forms are used throughout (\textit{quattuor, numquam, intellego, exspecto}).

Each page of this edition corresponds to one page of the 1610 edition; the running heads are those of the original, and the folio is given in brackets in the outer margin ([7r] is the recto of leaf 7). The Latin breaks exactly where the original does; the English breaks at the matching point.

The type follows the roles of the original: its roman as roman, its italic as italic, and words it sets in capitals as spaced capitals. The English translation is new and made directly from the Latin. The diagrams, star maps and drawings of Jupiter's companions are redrawn as vector graphics from measurements of the original; the engravings of the Moon, the woodcut initials, ornaments and printer's device are taken from the scan, cleaned and repaired.

The translation was made with AI assistance and has not yet been reviewed by a Latinist; corrections are welcome. Translation, notes and typesetting are released under CC BY 4.0; the Latin text and the scans of the 1610 edition are in the public domain. Sources: github.com/metachow/sidereus-nuncius
\end{minipage}
\cleardoublepage
